% sec10_2.tex — Section 10.2: Alternative Representations of Cointegrated Systems

In addition to the common trend representation (see Review Question~2 of the pre\-vious
section), there are three other useful representations of cointegrated vector
I(1) processes: the triangular representation of Phillips (1991), the VAR represen\-tation,
and the VECM (vector error-correction model) of Davidson et al.\ (1978).
This section introduces these representations.

\subsection*{Phillips's Triangular Representation}

This representation is convenient for the purpose of estimating cointegrating vec\-tors.
The representation is valid for any cointegration rank, but we initially assume
that the cointegrating rank $h$ is~1. Let $\mathbf{a}$ be a cointegrating vector, and suppose,
without loss of generality, that the first element of $\mathbf{a}$ is not zero (if it is zero, change
the ordering of the variables of the system so that the first element after reordering
is not zero). Partition $\mathbf{y}_t$ accordingly:
\begin{equation}
  \underset{(n \times 1)}{\mathbf{y}_t}
  \equiv \begin{bmatrix} y_{1t} \\
    \underset{((n-1) \times 1)}{\mathbf{y}_{2t}}
  \end{bmatrix}.
  \label{eq:10.2.1} \tag{10.2.1}
\end{equation}
Thus, $y_{1t}$ is cointegrated with $\mathbf{y}_{2t}$. Since a scalar multiple of a cointegrating vector,
too, is a cointegrating vector, we can normalize the cointegrating vector $\mathbf{a}$ so that
its first element is unity:
\begin{equation}
  \mathbf{a}
  = \begin{bmatrix} a_1 \\ a_2 \\ \vdots \\ a_n \end{bmatrix}
  = \begin{bmatrix} 1 \\
    \underset{((n-1) \times 1)}{-\boldsymbol{\gamma}}
  \end{bmatrix}.
  \label{eq:10.2.2} \tag{10.2.2}
\end{equation}

We have seen in the previous section that if a cointegrating vector $\mathbf{a}$ eliminates
not only the stochastic trend (i.e., $\mathbf{a}'\boldsymbol{\Psi}(1) = \mathbf{0}'$) but also the deterministic trend
(i.e., $\mathbf{a}'\boldsymbol{\delta} = 0$), then $\mathbf{a}'\mathbf{y}_t$ can be written as (10.1.24). Setting the $\mathbf{a}$ in (10.1.24) to
the cointegrating vector in (10.2.2), we obtain
\begin{equation}
  y_{1t} = \boldsymbol{\gamma}' \mathbf{y}_{2t} + z_t^* + \mu,
  \label{eq:10.2.3} \tag{10.2.3}
\end{equation}
where
\begin{equation}
  z_t^* \equiv (1, -\boldsymbol{\gamma}') \boldsymbol{\eta}_t, \quad
  \mu \equiv (1, -\boldsymbol{\gamma}')(\mathbf{y}_0 - \boldsymbol{\eta}_0).
  \label{eq:10.2.4} \tag{10.2.4}
\end{equation}
Since $\boldsymbol{\eta}_t$ is jointly stationary, $z_t^*$ is stationary. This equation, with $z_t^*$ viewed as an
error term and $\mu$ as an intercept, is called a \textbf{cointegrating regression}. Its regres\-sion
coefficients $\boldsymbol{\gamma}$, relating the permanent component in $y_{1t}$ to those in $\mathbf{y}_{2t}$, can
be interpreted as describing the long-run relationship between $y_{1t}$ and $\mathbf{y}_{2t}$. The
cointegrating regression will have the trend term $\mathbf{a}'\boldsymbol{\delta} \cdot t$ as an additional regressor
if the cointegrating vector does not eliminate the deterministic trend in $\mathbf{a}'\mathbf{y}_t$. The
\textbf{triangular representation} is an $n$-equation system consisting of this cointegrating
regression and the last $n - 1$ rows of (10.1.9):
\begin{equation}
  \underset{((n-1) \times 1)}{\Delta \mathbf{y}_{2t}}
  = \underset{((n-1) \times 1)}{\boldsymbol{\delta}_2}
    + \underset{((n-1) \times 1)}{\mathbf{u}_{2t}}
  = \underset{((n-1) \times 1)}{\boldsymbol{\delta}_2}
    + \underset{((n-1) \times n)}{\boldsymbol{\Psi}_2(L)} \;
      \underset{(n \times 1)}{\boldsymbol{\varepsilon}_t},
  \label{eq:10.2.5} \tag{10.2.5}
\end{equation}
where $\boldsymbol{\delta}_2$ and $\mathbf{u}_{2t}$ are the last $n - 1$ elements of the $n$-dimensional vectors $\boldsymbol{\delta}$ and
$\mathbf{u}_t$, respectively, and $\boldsymbol{\Psi}_2(L)$ is the last $n - 1$ rows of $\boldsymbol{\Psi}(L)$ in the VMA repre\-sentation
(10.1.10). The implication of the $h = 1$ assumption is that, as noted in
the previous section, $\mathbf{y}_{2t}$ is not cointegrated (it would be cointegrated if $h > 1$). In
particular, each element of $\mathbf{y}_{2t}$ is individually I(1).

More generally, consider the case where the cointegration rank $h$ is not neces\-sarily~1.
By changing the order of the elements of $\mathbf{y}_t$ if necessary, it is possible (see,
e.g., Hamilton, 1994, pp.~576--577) to select $h$ linearly independent cointegrating
vectors, $\mathbf{a}_1, \mathbf{a}_2, \ldots, \mathbf{a}_h$, such that
\begin{equation}
  \underset{(n \times h)}{\mathbf{A}}
  \equiv \begin{bmatrix} \mathbf{a}_1 & \cdots & \mathbf{a}_h \end{bmatrix}
  = \begin{bmatrix}
    \underset{(h \times h)}{\mathbf{I}_h} \\
    \underset{((n-h) \times h)}{-\boldsymbol{\Gamma}}
  \end{bmatrix}.
  \label{eq:10.2.6} \tag{10.2.6}
\end{equation}
Partition $\mathbf{y}_t$ conformably as
\begin{equation}
  \underset{(n \times 1)}{\mathbf{y}_t}
  \equiv \begin{bmatrix}
    \underset{(h \times 1)}{\mathbf{y}_{1t}} \\
    \underset{((n-h) \times 1)}{\mathbf{y}_{2t}}
  \end{bmatrix}.
  \label{eq:10.2.7} \tag{10.2.7}
\end{equation}
Multiplying both sides of the BN decomposition (10.1.16) by this $\mathbf{A}'$ and noting
that $\mathbf{A}'\boldsymbol{\Psi}(1) = \mathbf{0}$ (since the columns of $\mathbf{A}$ are cointegrating vectors),
$\mathbf{A}'\boldsymbol{\delta} = \mathbf{0}$
(if those cointegrating vectors also eliminate the deterministic trend), and
$\mathbf{A}'\mathbf{y}_t = \mathbf{y}_{1t} - \boldsymbol{\Gamma}'\mathbf{y}_{2t}$, we obtain the following $h$ cointegrating regressions:
\begin{equation}
  \underset{(h \times 1)}{\mathbf{y}_{1t}}
  = \underset{(h \times (n-h))}{\boldsymbol{\Gamma}'} \;
    \underset{((n-h) \times 1)}{\mathbf{y}_{2t}}
    + \underset{(h \times 1)}{\boldsymbol{\mu}}
    + \underset{(h \times 1)}{\mathbf{z}_t^*},
  \label{eq:10.2.8} \tag{10.2.8}
\end{equation}
where $\boldsymbol{\mu} \equiv \mathbf{A}'(\mathbf{y}_0 - \boldsymbol{\eta}_0)$ and
$\mathbf{z}_t^* \equiv \mathbf{A}'\boldsymbol{\eta}_t$. Since $\boldsymbol{\eta}_t$ is jointly stationary, so is $\mathbf{z}_t^*$. The
triangular representation is these $h$ cointegrating regressions, supplemented by the
rest of the VMA representation:
\begin{equation}
  \underset{((n-h) \times 1)}{\Delta \mathbf{y}_{2t}}
  = \underset{((n-h) \times 1)}{\boldsymbol{\delta}_2}
    + \underset{((n-h) \times 1)}{\mathbf{u}_{2t}}
  = \underset{((n-h) \times 1)}{\boldsymbol{\delta}_2}
    + \underset{((n-h) \times n)}{\boldsymbol{\Psi}_2(L)} \;
      \underset{(n \times 1)}{\boldsymbol{\varepsilon}_t}.
  \label{eq:10.2.9} \tag{10.2.9}
\end{equation}
Here, $\boldsymbol{\Psi}_2(L)$ is the last $n - h$ rows of the $\boldsymbol{\Psi}(L)$ in the VMA representation (10.1.10).
It is easy to show (see Review Question~1) that $\mathbf{y}_{2t}$ is not cointegrated.

To illustrate the triangular representation and how it can be derived from the
VMA representation, consider

\begin{defbox}
\textbf{Example 10.3\quad (From VMA to triangular):}\quad In the bivariate system (10.1.12)
of Example~10.2, the cointegration rank is~1. The cointegrating vector whose
first element is unity is $(1, -\gamma)'$. So the $z_t^*$ and $\mu$ in (10.2.3) are
\begin{align}
  z_t^* &= (1, -\gamma)\boldsymbol{\eta}_t = \varepsilon_{1t} - \gamma \varepsilon_{2t},
  \notag \\
  \mu &= (1, -\gamma)(\mathbf{y}_0 - \boldsymbol{\eta}_0)
    = (y_{1,0} - \gamma y_{2,0}) - (\varepsilon_{1,0} - \gamma \varepsilon_{2,0}),
  \label{eq:10.2.10} \tag{10.2.10}
\end{align}
and the triangular representation is
\begin{equation}
  \begin{cases}
    y_{1t} = \mu + \gamma y_{2t} + (\varepsilon_{1t} - \gamma \varepsilon_{2t}), \\
    \Delta y_{2t} = \delta_2 + \varepsilon_{2t}.
  \end{cases}
  \label{eq:10.2.11} \tag{10.2.11}
\end{equation}
\end{defbox}

\subsection*{VAR and Cointegration}

For the stationary case, we found it useful and convenient to model a vector process
as a finite-order VAR. Although, as seen above, no cointegrated I(1) system can be
represented as a finite-order VAR in first differences, some cointegrated systems
may admit a finite-order VAR representation in \emph{levels}. So suppose a cointegrated
I(1) system $\mathbf{y}_t$ can be written as
\begin{align}
  \underset{(n \times 1)}{\mathbf{y}_t}
  &= \boldsymbol{\alpha} + \boldsymbol{\delta} \cdot t + \boldsymbol{\xi}_t, \quad
  \boldsymbol{\Phi}(L) \boldsymbol{\xi}_t = \boldsymbol{\varepsilon}_t,
  \notag \\
  \underset{(n \times n)}{\boldsymbol{\Phi}(L)}
  &\equiv \mathbf{I}_n - \boldsymbol{\Phi}_1 L - \boldsymbol{\Phi}_2 L^2
    - \cdots - \boldsymbol{\Phi}_p L^p.
  \label{eq:10.2.12} \tag{10.2.12}
\end{align}
For later reference, we eliminate $\boldsymbol{\xi}$ from this to obtain
\begin{equation}
  \mathbf{y}_t = \boldsymbol{\alpha}^* + \boldsymbol{\delta}^* \cdot t
    + \boldsymbol{\Phi}_1 \mathbf{y}_{t-1} + \boldsymbol{\Phi}_2 \mathbf{y}_{t-2}
    + \cdots + \boldsymbol{\Phi}_p \mathbf{y}_{t-p} + \boldsymbol{\varepsilon}_t,
  \label{eq:10.2.13} \tag{10.2.13}
\end{equation}
where
\begin{equation}
  \underset{(n \times 1)}{\boldsymbol{\alpha}^*}
  \equiv (\boldsymbol{\Phi}_1 + 2\boldsymbol{\Phi}_2 + \cdots + p\boldsymbol{\Phi}_p)\boldsymbol{\delta}
    + \boldsymbol{\Phi}(1)\boldsymbol{\alpha}, \quad
  \underset{(n \times 1)}{\boldsymbol{\delta}^*}
  \equiv \boldsymbol{\Phi}(1)\boldsymbol{\delta}.
  \label{eq:10.2.14} \tag{10.2.14}
\end{equation}

How do we know that this finite-order VAR in levels is a cointegrated I(1)
system? An obvious way to find out is to derive the VMA representation,
$\Delta\boldsymbol{\xi}_t = \boldsymbol{\Psi}(L)\boldsymbol{\varepsilon}_t$,
from the VAR representation, $\boldsymbol{\Phi}(L)\boldsymbol{\xi}_t = \boldsymbol{\varepsilon}_t$, and see if $\boldsymbol{\Psi}(L)$ satisfies the
definition of a cointegrated system. The derivation is a bit tricky because the VMA
representation is in first differences, but it can be done fairly easily as follows.
Taking the first difference of both sides of $\boldsymbol{\Phi}(L)\boldsymbol{\xi}_t = \boldsymbol{\varepsilon}_t$ and noting that
$\Delta = 1 - L$ and $(1 - L)\boldsymbol{\Phi}(L) = \boldsymbol{\Phi}(L)(1 - L)$, we obtain
\begin{equation}
  \boldsymbol{\Phi}(L) \Delta\boldsymbol{\xi}_t = (1 - L)\boldsymbol{\varepsilon}_t.
  \label{eq:10.2.15} \tag{10.2.15}
\end{equation}
Substitute $\Delta\boldsymbol{\xi}_t = \boldsymbol{\Psi}(L)\boldsymbol{\varepsilon}_t$ into this to obtain
\begin{equation}
  \boldsymbol{\Phi}(L) \boldsymbol{\Psi}(L) \boldsymbol{\varepsilon}_t
  = (1 - L)\boldsymbol{\varepsilon}_t.
  \label{eq:10.2.16} \tag{10.2.16}
\end{equation}
This equation has to hold for any realization of $\boldsymbol{\varepsilon}_t$, so
\begin{equation}
  \boldsymbol{\Phi}(L) \boldsymbol{\Psi}(L) = (1 - L)\mathbf{I}_n.
  \label{eq:10.2.17} \tag{10.2.17}
\end{equation}
This can be solved for $\boldsymbol{\Psi}(L)$ by multiplying both sides from the left by
$\boldsymbol{\Phi}(L)^{-1}$, the
inverse of $\boldsymbol{\Phi}(L)$.\footnote{%
The inverse exists because $\boldsymbol{\Phi}_0 = \mathbf{I}_n$ is nonsingular; see Section~6.3. Since we are not assuming the station\-arity
condition for $\boldsymbol{\Phi}(L)$, the inverse filter may not be absolutely summable.}
This produces: $\boldsymbol{\Psi}(L) = \boldsymbol{\Phi}(L)^{-1}(1 - L)$. The question is, under
what conditions on $\boldsymbol{\Phi}(L)$ is this $\boldsymbol{\Psi}(L)$ one-summable and
$\rank[\boldsymbol{\Psi}(1)] = n - h$?

We can easily derive a necessary condition. Setting $L = 1$ in (10.2.17), we
obtain
\begin{equation}
  \underset{(n \times n)}{\boldsymbol{\Phi}(1)} \;
  \underset{(n \times n)}{\boldsymbol{\Psi}(1)}
  = \underset{(n \times n)}{\mathbf{0}}.
  \label{eq:10.2.18} \tag{10.2.18}
\end{equation}
Since the rank of $\boldsymbol{\Psi}(1)$ equals $n - h$ when the cointegration rank of $\boldsymbol{\xi}_t$ is $h$, the rank
of $\boldsymbol{\Phi}(1)$ is at most $h$. As shown below, the rank is actually $h$. To state a necessary
and sufficient condition, let $\mathbf{U}(L)$ and $\mathbf{V}(L)$ be $n \times n$ matrix lag polynomials with
all their roots outside the unit circle and let $\mathbf{M}(L)$ be a matrix lag polynomial satisfying
\begin{equation}
  \mathbf{M}(L)
  = \begin{bmatrix}
    (1 - L)\mathbf{I}_{n-h} & \underset{((n-h) \times h)}{\mathbf{0}} \\
    \underset{(h \times (n-h))}{\mathbf{0}} & \mathbf{I}_h
  \end{bmatrix}.
  \label{eq:10.2.19} \tag{10.2.19}
\end{equation}
That is, the first $n - h$ diagonal elements of the diagonal matrix $\mathbf{M}(L)$ are $1 - L$ and
the remaining $h$ diagonal elements are unity. A necessary and sufficient condition
for a finite-order VAR process $\{\boldsymbol{\xi}_t\}$ following $\boldsymbol{\Phi}(L)\boldsymbol{\xi}_t = \boldsymbol{\varepsilon}_t$ to be a cointegrated
I(1) system with rank $h$ is that $\boldsymbol{\Phi}(L)$ can be factored as
$\boldsymbol{\Phi}(L) = \mathbf{U}(L)\mathbf{M}(L)\mathbf{V}(L)$.\footnote{%
This result is an implication of the lemma due to Sam Yoo, cited in Engle and Yoo (1991). See Watson (1994,
pp.~2870--2873) for an accessible exposition.}
Therefore, all the roots of $|\boldsymbol{\Phi}(z)| = 0$ are on or outside the unit circle and those
that are on the unit circle are all unit roots ($z = 1$). It is not sufficient that $\boldsymbol{\Phi}(L)$
has $n - h$ unit roots with the other $h$ roots outside the unit circle; see an example
in Review Question~4 where $\boldsymbol{\Phi}(z)$ has two unit roots and one root outside the unit
circle yet the system is not I(1). The $n - h$ unit roots have to be located in the
system in the particular way indicated by the factorization
$\boldsymbol{\Phi}(z) = \mathbf{U}(z)\mathbf{M}(z)\mathbf{V}(z)$.

Setting $z = 1$ in this factorization, we obtain
$\boldsymbol{\Phi}(1) = \mathbf{U}(1)\mathbf{M}(1)\mathbf{V}(1)$. Since the
roots of $\mathbf{U}(z)$ and $\mathbf{V}(z)$ are all outside the unit circle, $\mathbf{U}(1)$ and $\mathbf{V}(1)$ are nonsingular
and the rank of $\boldsymbol{\Phi}(1)$ equals the rank of $\mathbf{M}(1)$, which is $h$. That is,
\begin{equation}
  \rank[\boldsymbol{\Phi}(1)] = h.
  \label{eq:10.2.20} \tag{10.2.20}
\end{equation}
Then it follows from basic linear algebra that there exist two $n \times h$ matrices of full
column rank, $\mathbf{B}$ and $\mathbf{A}$, such that
\begin{equation}
  \underset{(n \times n)}{\boldsymbol{\Phi}(1)}
  = \underset{(n \times h)}{\mathbf{B}} \;
    \underset{(h \times n)}{\mathbf{A}'}.
  \label{eq:10.2.21} \tag{10.2.21}
\end{equation}
This is sometimes called the \textbf{reduced rank condition}. The choice of $\mathbf{B}$ and $\mathbf{A}$ is
not unique; if $\mathbf{F}$ is an $h \times h$ nonsingular matrix, then $\mathbf{B}(\mathbf{F}')^{-1}$ and $\mathbf{AF}$ in place
of $\mathbf{B}$ and $\mathbf{A}$ also satisfy (10.2.21). Substituting (10.2.21) into (10.2.18), we obtain
$\mathbf{B}\mathbf{A}'\boldsymbol{\Psi}(1) = \mathbf{0}$. Since $\mathbf{B}$ is of full column rank, this equation implies
$\mathbf{A}'\boldsymbol{\Psi}(1) = \mathbf{0}$.
So the columns of the $n \times h$ matrix $\mathbf{A}$ are cointegrating vectors.

\subsection*{The Vector Error-Correction Model (VECM)}

For the univariate case, we derived the augmented autoregression (9.4.3) on page
586 from an AR equation. The same idea can be applied to the VAR here. It is
a matter of simple algebra to show that the VAR representation $\boldsymbol{\Phi}(L)\boldsymbol{\xi}_t = \boldsymbol{\varepsilon}_t$ in
(10.2.12) can be written equivalently as
\begin{equation}
  \boldsymbol{\xi}_t = \boldsymbol{\rho}\,\boldsymbol{\xi}_{t-1}
    + \boldsymbol{\zeta}_1 \Delta\boldsymbol{\xi}_{t-1}
    + \boldsymbol{\zeta}_2 \Delta\boldsymbol{\xi}_{t-2}
    + \cdots
    + \boldsymbol{\zeta}_{p-1} \Delta\boldsymbol{\xi}_{t-p+1}
    + \boldsymbol{\varepsilon}_t,
  \label{eq:10.2.22} \tag{10.2.22}
\end{equation}
where
\begin{align}
  \underset{(n \times n)}{\boldsymbol{\rho}}
  &\equiv \boldsymbol{\Phi}_1 + \boldsymbol{\Phi}_2 + \cdots + \boldsymbol{\Phi}_p,
  \label{eq:10.2.23} \tag{10.2.23} \\
  \underset{(n \times n)}{\boldsymbol{\zeta}_s}
  &\equiv -(\boldsymbol{\Phi}_{s+1} + \boldsymbol{\Phi}_{s+2} + \cdots + \boldsymbol{\Phi}_p)
  \quad \text{for } s = 1, 2, \ldots, p - 1.
  \label{eq:10.2.24} \tag{10.2.24}
\end{align}
Subtracting $\boldsymbol{\xi}_{t-1}$ from both sides of (10.2.22) and noting that
$\boldsymbol{\rho} - \mathbf{I}_n = -(\mathbf{I}_n - \boldsymbol{\Phi}_1 - \boldsymbol{\Phi}_2 - \cdots - \boldsymbol{\Phi}_p) = -\boldsymbol{\Phi}(1)$,
we can rewrite (10.2.22) as
\begin{align}
  \Delta\boldsymbol{\xi}_t
  &= -\boldsymbol{\Phi}(1)\boldsymbol{\xi}_{t-1}
    + \boldsymbol{\zeta}_1 \Delta\boldsymbol{\xi}_{t-1}
    + \boldsymbol{\zeta}_2 \Delta\boldsymbol{\xi}_{t-2}
    + \cdots
    + \boldsymbol{\zeta}_{p-1} \Delta\boldsymbol{\xi}_{t-p+1}
    + \boldsymbol{\varepsilon}_t
  \notag \\
  &= -\mathbf{B}\mathbf{A}'\boldsymbol{\xi}_{t-1}
    + \boldsymbol{\zeta}_1 \Delta\boldsymbol{\xi}_{t-1}
    + \boldsymbol{\zeta}_2 \Delta\boldsymbol{\xi}_{t-2}
    + \cdots
    + \boldsymbol{\zeta}_{p-1} \Delta\boldsymbol{\xi}_{t-p+1}
    + \boldsymbol{\varepsilon}_t.
  \label{eq:10.2.25} \tag{10.2.25}
\end{align}
Using the relation $\mathbf{y}_t = \boldsymbol{\alpha} + \boldsymbol{\delta} \cdot t + \boldsymbol{\xi}_t$, this can be translated into an equation in $\mathbf{y}_t$:
\begin{align}
  \Delta \mathbf{y}_t
  &= \boldsymbol{\alpha}^* + \boldsymbol{\delta}^* \cdot t
    - \boldsymbol{\Phi}(1)\,\mathbf{y}_{t-1}
    + \boldsymbol{\zeta}_1 \Delta\mathbf{y}_{t-1}
    + \cdots
    + \boldsymbol{\zeta}_{p-1} \Delta\mathbf{y}_{t-p+1}
    + \boldsymbol{\varepsilon}_t
  \label{eq:10.2.26} \tag{10.2.26} \\
  &= \boldsymbol{\alpha}^* + \boldsymbol{\delta}^* \cdot t
    - \underset{(n \times h)}{\mathbf{B}} \;
      \underset{(h \times 1)}{\mathbf{z}_{t-1}}
    + \boldsymbol{\zeta}_1 \Delta\mathbf{y}_{t-1}
    + \cdots
    + \boldsymbol{\zeta}_{p-1} \Delta\mathbf{y}_{t-p+1}
    + \boldsymbol{\varepsilon}_t,
  \label{eq:10.2.27} \tag{10.2.27}
\end{align}
where $\boldsymbol{\alpha}^*$ and $\boldsymbol{\delta}^*$ are as in (10.2.14) and $\mathbf{z}_t$ is given by
\begin{equation}
  \underset{(h \times 1)}{\mathbf{z}_t}
  \equiv \underset{(h \times n)}{\mathbf{A}'} \;
    \underset{(n \times 1)}{\mathbf{y}_t}.
  \label{eq:10.2.28} \tag{10.2.28}
\end{equation}
Since, as noted above, the $n \times h$ matrix $\mathbf{A}$ collects $h$ cointegrating vectors, $\mathbf{z}_t$ is
trend stationary (with a suitable choice of the initial value $\mathbf{y}_0$).\footnote{%
Just in case you are wondering why $\mathbf{y}_0$ is relevant. It is true that you can solve (10.2.27) for $\mathbf{z}_t$ as
$\mathbf{z}_{t-1} = (\mathbf{B}'\mathbf{B})^{-1}\mathbf{B}'[\boldsymbol{\alpha}^* + \boldsymbol{\delta}^* \cdot t
  - \Delta \mathbf{y}_t + \boldsymbol{\zeta}_1 \Delta\mathbf{y}_{t-1} + \cdots + \boldsymbol{\zeta}_{p-1}\Delta\mathbf{y}_{t-p+1} + \boldsymbol{\varepsilon}_t]$.
So you might think that $\mathbf{z}_t$ is trend stationary regardless of the choice of $\mathbf{y}_0$. This is not true, strictly speak\-ing,
because, unlike the VMA representation, $\Delta\mathbf{y}_t$ in the VAR/VECM representation is not defined for
$t = -1, -2, \ldots$\,. Only with a judicious choice of $\mathbf{y}_0$ can one make $\Delta\mathbf{y}_t$ ($t = 1, 2, \ldots$) stationary. This point is
mentioned in Johansen (1995, Theorem~4.2).}
This representa\-tion
is called the \textbf{vector error-correction model (VECM)} or the \textbf{error-correction
representation}. If it were not for the term $\mathbf{B}\mathbf{z}_{t-1}$ in the VECM, the process,
expressed as a VAR in first differences, could not be cointegrated. The VECM
accommodates $h$ cointegrating relationships by including $h$ linear combinations of
levels of the variables in the cointegrating relations. If there are no time trends in the
cointegrating relations,
that is, if $\mathbf{A}'\boldsymbol{\delta} = \mathbf{0}$, then $\boldsymbol{\delta}^* = \boldsymbol{\Phi}(1)\boldsymbol{\delta} = \mathbf{B}\mathbf{A}'\boldsymbol{\delta} = \mathbf{0}$. So the VAR and the VECM
representations do not involve time trends despite the possible existence of time
trends in the elements of $\mathbf{y}_t$.

That the same I(0) process has the VMA, VAR, and VECM representations is
known as the \textbf{Granger Representation Theorem}.

\begin{defbox}
\textbf{Example 10.4\quad (From VMA to VAR/VECM):}\quad In the previous example, we
derived the triangular representation from the VMA representation (10.1.12).
In this example, we derive the VAR and VECM representations from the
same VMA representation. For the $\boldsymbol{\Psi}(L)$ in (10.1.12), it is easy to verify that
(10.2.17) is satisfied with
\begin{equation}
  \boldsymbol{\Phi}(L) = \mathbf{I}_2 - \boldsymbol{\Phi}_1 L, \quad
  \boldsymbol{\Phi}_1
  = \begin{bmatrix} 0 & \gamma \\ 0 & 1 \end{bmatrix}.
  \label{eq:10.2.29} \tag{10.2.29}
\end{equation}
So the VMA can be represented by a finite-order VAR. For this VAR, we have
\begin{equation}
  \boldsymbol{\Phi}(1) = \mathbf{I}_2 - \boldsymbol{\Phi}_1
  = \begin{bmatrix} 1 & -\gamma \\ 0 & 0 \end{bmatrix},
  \label{eq:10.2.30} \tag{10.2.30}
\end{equation}
which can be written as (10.2.21) with
\begin{equation}
  \mathbf{B} = \begin{bmatrix} 1 \\ 0 \end{bmatrix}
  \quad \text{and} \quad
  \mathbf{A} = \begin{bmatrix} 1 \\ -\gamma \end{bmatrix}
  \quad \text{(for example)}.
  \label{eq:10.2.31} \tag{10.2.31}
\end{equation}
By (10.2.27) and (10.2.28), the VECM for this choice of $\mathbf{B}$ and $\mathbf{A}$ is
\begin{equation}
  \begin{bmatrix} \Delta y_{1t} \\ \Delta y_{2t} \end{bmatrix}
  = \begin{bmatrix} \gamma\delta_2 + \alpha_1 - \gamma\alpha_2 \\
                     \delta_2 \end{bmatrix}
  + \begin{bmatrix} \delta_1 - \gamma\delta_2 \\ 0 \end{bmatrix} \cdot t
  - \begin{bmatrix} 1 \\ 0 \end{bmatrix} z_{t-1}
  + \boldsymbol{\varepsilon}_t,
  \quad z_t \equiv \mathbf{A}' \mathbf{y}_t = y_{1t} - \gamma y_{2t}.
  \label{eq:10.2.32} \tag{10.2.32}
\end{equation}
The deterministic trend disappears if $\delta_1 = \gamma\delta_2$, that is, if the cointe\-grating
vector also eliminates the deterministic trend.
\end{defbox}

\subsection*{Johansen's ML Procedure}

In closing this section, having introduced the VECM representation, we here
provide an executive summary of the maximum likelihood (ML) estimation of
cointegrated systems proposed by Johansen (1988). Go back to the VAR repre\-sentation
(10.2.13). For simplicity, we assume that there is no trend in the
cointegrating relations, so that $\boldsymbol{\delta}^* = \mathbf{0}$.\footnote{%
For a thorough treatment of time trends in the VECM, see Johansen (1995, Sections~5.7 and~6.2).}
We have considered the conditional ML
estimation (conditional on the initial values of $\mathbf{y}$) of a VAR in Section~8.7. If the
error vector $\boldsymbol{\varepsilon}_t$ is jointly normal $N(\mathbf{0}, \boldsymbol{\Omega})$ and if we have a sample of $T + p$ obser\-vations
$(\mathbf{y}_{-p+1}, \mathbf{y}_{-p+2}, \ldots, \mathbf{y}_T)$, then the (average) log likelihood of
$(\mathbf{y}_1, \ldots, \mathbf{y}_T)$
conditional on the initial values $(\mathbf{y}_{-p+1}, \mathbf{y}_{-p+2}, \ldots, \mathbf{y}_0)$ is given by
\begin{equation}
  Q_T(\boldsymbol{\theta})
  = -\frac{n}{2}\log(2\pi)
    + \frac{1}{2}\log(|\boldsymbol{\Omega}^{-1}|)
    - \frac{1}{2T} \sum_{t=1}^{T}
      \bigl\{(\mathbf{y}_t - \boldsymbol{\Pi}'\mathbf{x}_t)'
        \boldsymbol{\Omega}^{-1}
        (\mathbf{y}_t - \boldsymbol{\Pi}'\mathbf{x}_t)\bigr\},
  \label{eq:10.2.33} \tag{10.2.33}
\end{equation}
where $n$ is the dimension of the system (not the sample size),
$\boldsymbol{\theta} = (\boldsymbol{\Pi}, \boldsymbol{\Omega})$, and
\begin{equation}
  \boldsymbol{\Pi}' \equiv \begin{bmatrix}
    \underset{(n \times 1)}{\boldsymbol{\alpha}^*} &
    \underset{(n \times n)}{\boldsymbol{\Phi}_1} &
    \underset{(n \times n)}{\boldsymbol{\Phi}_2} &
    \cdots &
    \underset{(n \times n)}{\boldsymbol{\Phi}_p}
  \end{bmatrix}, \quad
  \mathbf{x}_t \equiv \begin{bmatrix}
    1 \\ \mathbf{y}_{t-1} \\ \mathbf{y}_{t-2} \\ \vdots \\ \mathbf{y}_{t-p}
  \end{bmatrix}.
  \label{eq:10.2.34} \tag{10.2.34}
\end{equation}

Let
\begin{equation}
  \underset{(n \times n)}{\boldsymbol{\zeta}_0}
  \equiv -\boldsymbol{\Phi}(1)
  = -\bigl[\mathbf{I}_n - (\boldsymbol{\Phi}_1 + \cdots + \boldsymbol{\Phi}_p)\bigr].
  \label{eq:10.2.35} \tag{10.2.35}
\end{equation}
Then the reduced rank condition is that $\boldsymbol{\zeta}_0 = -\mathbf{B}\mathbf{A}'$ and the VECM (10.2.26) can
be written as
\begin{equation}
  \Delta \mathbf{y}_t = \boldsymbol{\alpha}^*
    + \boldsymbol{\zeta}_0 \mathbf{y}_{t-1}
    + \boldsymbol{\zeta}_1 \Delta\mathbf{y}_{t-1}
    + \cdots
    + \boldsymbol{\zeta}_{p-1} \Delta\mathbf{y}_{t-p+1}
    + \boldsymbol{\varepsilon}_t.
  \label{eq:10.2.36} \tag{10.2.36}
\end{equation}
Equations (10.2.35) and (10.2.24) provide a one-to-one mapping between
$(\boldsymbol{\Phi}_1, \ldots, \boldsymbol{\Phi}_p)$ and
$(\boldsymbol{\zeta}_0, \ldots, \boldsymbol{\zeta}_{p-1})$. Thus the same average log likelihood $Q_T(\boldsymbol{\theta})$ can be
rewritten in terms of the VECM parameters as
\begin{equation}
  -\frac{n}{2}\log(2\pi)
  + \frac{1}{2}\log(|\boldsymbol{\Omega}^{-1}|)
  - \frac{1}{2T} \sum_{t=1}^{T}
    \bigl\{(\Delta\mathbf{y}_t - \tilde{\boldsymbol{\Pi}}'\tilde{\mathbf{x}}_t)'
      \boldsymbol{\Omega}^{-1}
      (\Delta\mathbf{y}_t - \tilde{\boldsymbol{\Pi}}'\tilde{\mathbf{x}}_t)\bigr\},
  \label{eq:10.2.37} \tag{10.2.37}
\end{equation}
where
\begin{equation}
  \tilde{\boldsymbol{\Pi}}' \equiv \begin{bmatrix}
    \underset{(n \times 1)}{\boldsymbol{\alpha}^*} &
    \underset{(n \times n)}{\boldsymbol{\zeta}_0} &
    \underset{(n \times n)}{\boldsymbol{\zeta}_1} &
    \cdots &
    \underset{(n \times n)}{\boldsymbol{\zeta}_{p-1}}
  \end{bmatrix}, \quad
  \tilde{\mathbf{x}}_t \equiv \begin{bmatrix}
    1 \\ \mathbf{y}_{t-1} \\ \Delta\mathbf{y}_{t-1} \\ \vdots \\ \Delta\mathbf{y}_{t-p+1}
  \end{bmatrix}.
  \label{eq:10.2.38} \tag{10.2.38}
\end{equation}
The ML estimate of the VECM parameters is the
$(\boldsymbol{\alpha}^*, \boldsymbol{\zeta}_0, \ldots, \boldsymbol{\zeta}_{p-1}, \boldsymbol{\Omega})$ that max\-imizes
this objective function, subject to the reduced rank constraint that $\boldsymbol{\zeta}_0 =
-\mathbf{B}\mathbf{A}'$ for some $n \times h$ full rank matrices $\mathbf{A}$ and $\mathbf{B}$. This constraint accommodates
$h$ cointegrating relationships on the I(1) system. Obviously, the maximized log
likelihood is higher the higher the assumed cointegration rank $h$. Thus, we can use
the likelihood ratio statistic to test the null hypothesis that $h = h_0$ against the alter\-native
hypothesis of \emph{more} cointegration. Unlike in the stationary VAR case, the
limiting distribution of the likelihood ratio statistic is nonstandard. Given the coin\-tegration
rank $h$ thus determined, the ML estimate of $h$ cointegrating vectors can
be obtained as the estimate of $\mathbf{A}$. For a more detailed exposition of this procedure,
see Hamilton (1994, Chapter~20) and also Johansen (1995, Chapter~6).

\subsection*{Questions for Review}

\begin{enumerate}
\item (Error vector in the triangular representation)\quad Consider the error vector
  $(\mathbf{z}_t^*, \mathbf{u}_{2t})$ in the triangular representation (10.2.8) and (10.2.9). It can be writ\-ten
  as
  \begin{equation}
    \begin{bmatrix}
      \underset{(h \times 1)}{\mathbf{z}_t^*} \\
      \underset{((n-h) \times 1)}{\mathbf{u}_{2t}}
    \end{bmatrix}
    = \boldsymbol{\Psi}^*(L) \boldsymbol{\varepsilon}_t
    \equiv \begin{bmatrix}
      \underset{(h \times n)}{\boldsymbol{\Psi}_1^*(L)} \\
      \underset{((n-h) \times n)}{\boldsymbol{\Psi}_2(L)}
    \end{bmatrix} \boldsymbol{\varepsilon}_t.
    \label{eq:10.2.39} \tag{10.2.39}
  \end{equation}
  Here,
  \begin{equation}
    \underset{(h \times n)}{\boldsymbol{\Psi}_1^*(L)}
    = (\mathbf{I}, -\boldsymbol{\Gamma}') \;
      \underset{(n \times n)}{\boldsymbol{\alpha}(L)},
    \label{eq:10.2.40} \tag{10.2.40}
  \end{equation}
  where $\boldsymbol{\alpha}(L)$ is from the BN decomposition and $\boldsymbol{\Psi}_2(L)$ is the last $n - h$ rows of
  $\boldsymbol{\Psi}(L)$ in the VMA representation.

  \begin{enumerate}[label=(\alph*)]
  \item Show that $\boldsymbol{\Psi}_2(1)$ is of full row rank (i.e., the $n - h$ rows of $\boldsymbol{\Psi}_2(1)$ are linearly
    independent). \textbf{Hint:} Suppose, contrary to the claim, that there exists an
    $(n-h)$-dimensional vector $\mathbf{b} \neq \mathbf{0}$ such that
    $\mathbf{b}'\boldsymbol{\Psi}_2(1) = \mathbf{0}'$. Show that $(\mathbf{0}', \mathbf{b}')$
    would be an ($n$-dimensional) cointegrating vector and the cointegration rank
    would be at least $h + 1$.

  \item Verify that $\boldsymbol{\Psi}^*(L)$ is absolutely summable.

  \item Write $\boldsymbol{\Psi}^*(L) = \boldsymbol{\Psi}_0^* + \boldsymbol{\Psi}_1^* L + \boldsymbol{\Psi}_2^* L^2 + \cdots$.
    For the bivariate process of
    Example~10.3, verify that $\boldsymbol{\Psi}_0^*$ is not diagonal. (Therefore, even if the ele\-ments
    of $\boldsymbol{\varepsilon}_t$ are uncorrelated, $\mathbf{z}_t^*$ and $\mathbf{u}_{2t}$ can be correlated.)
  \end{enumerate}

\item (From triangular to VMA representations)\quad For Example~10.3, start from the
  triangular representation (10.2.11) and recover the VMA representation. \textbf{Hint:}
  Take the first difference of the cointegrating regression.

\item (An alternative decomposition of $\boldsymbol{\Phi}(1)$)\quad For the bivariate system of Example
  10.4, verify that
  \[
    \mathbf{B} = (\gamma, 0)', \quad \mathbf{A} = (1/\gamma, -1)'
  \]
  is an alternative decomposition of $\boldsymbol{\Phi}(1)$. Write down the VECM corresponding
  to this choice of $\mathbf{B}$ and $\mathbf{A}$.

\item (A trivariate VAR that is I(2))\quad Consider a trivariate VAR given by
  \[
    \begin{cases}
      y_{1t} = 2y_{1,t-1} - y_{1,t-2} + \varepsilon_{1t}
        \quad (\text{i.e., } \Delta^2 y_{1t} = \varepsilon_{1t}), \\
      y_{2t} = \rho\, y_{2,t-1} + \varepsilon_{2t}
        \quad \text{with } |\rho| < 1, \text{ and} \\
      y_{3t} = \varepsilon_{3t}.
    \end{cases}
  \]
  Verify that this is an I(2) system by writing $\Delta^2 \mathbf{y}_t$ as a vector moving average
  process. Write down $\boldsymbol{\Phi}(L)$ for this system and show that $\boldsymbol{\Phi}(z)$ has three roots,
  two unit roots and one that is outside the unit circle.
\end{enumerate}
